\chapter{Reference Guide: Public API}
\label{app:reference}

This appendix consolidates the engine's public API for use as a standing
reference once the narrative chapters of Parts~I--III have been read. It
mirrors the project's own \texttt{reference\_guide.md} in full, and
additionally interleaves short discussion paragraphs connecting each API
surface back to the chapter in which its design rationale was developed, so
that this appendix can serve a reader either as a pure lookup table or as a
condensed second pass through the book's argument.

\section{Import Conventions}

\begin{lstlisting}
from dagon import Workflow, DataMover, StagerMover
from dagon.task import DagonTask, TaskType
\end{lstlisting}

These two import lines are, in a precise sense, the entire public surface a
workflow author needs for the great majority of workflows discussed in this
book: every example from \Cref{ch:lessons-01-02} through
\Cref{ch:native-tasks} is expressible using only the names imported here,
plus the task-type-specific keyword arguments documented later in this
appendix. The optional \texttt{dagon.api}, \texttt{dagon.dag\_tps}, and
\texttt{dagon.config} modules, discussed in their respective sections below,
are imported separately precisely because they are optional in the sense
already developed in \Cref{ch:getting-started}: a workflow that never uses
the DAGon service, meta-workflows, or programmatic configuration access has
no need to import them at all.

\section{\texttt{dagon.Workflow}}

\begin{lstlisting}[caption={Constructor signature.}]
Workflow(
    name,
    config=None,
    config_file="dagon.ini",
    max_threads=10,
    jsonload=None,
    checkpoint_file=None,
)
\end{lstlisting}

\begin{description}[leftmargin=2.2em, style=nextline]
  \item[\texttt{name}] workflow name.
  \item[\texttt{config}] optional in-memory configuration dictionary.
  \item[\texttt{config\_file}] INI file path used when \texttt{config} is
    not provided.
  \item[\texttt{max\_threads}] semaphore size currently stored for task use.
  \item[\texttt{jsonload}] JSON-like workflow object to load.
  \item[\texttt{checkpoint\_file}] file path for writing checkpoint
    metadata.
\end{description}

Recall from \Cref{ch:configuration}'s discussion of the layered override
mechanism that \texttt{config} and \texttt{config\_file} are mutually
exclusive in effect, not additive: supplying \texttt{config} bypasses
\texttt{config\_file} entirely, rather than overlaying it. The
\texttt{checkpoint\_file} parameter, examined in full in
\Cref{ch:lessons-05-07} and Appendix~D, is what activates the
\texttt{<workflow-name>.<task-name>} checkpoint-key scheme; omitting it
entirely disables checkpoint recording, which is the correct choice for the
many short, cheap-to-rerun tutorial-scale workflows used throughout
Part~I, and the incorrect choice for any workflow expensive enough to
warrant resumability.

\paragraph{Selected methods.}
\begin{itemize}[leftmargin=1.4em]
  \item \texttt{add\_task(task)}: add a task and attach it to this workflow.
  \item \texttt{make\_dependencies()}: parse dataflow references and
    validate the graph.
  \item \texttt{run(resume\_checkpoint\_file=None)}: run tasks in the
    current thread.
  \item \texttt{launch(resume\_checkpoint\_file=None)}: run tasks in a
    background thread.
  \item \texttt{wait(timeout=None)}: wait for a launched workflow; returns
    whether it ended.
  \item \texttt{add\_listener(event\_name, callback)}: register a workflow
    or task lifecycle callback.
  \item \texttt{set\_dry(dry)} / \texttt{get\_dry()}: set/return the
    dry-run flag.
  \item \texttt{set\_data\_mover(data\_mover)} / \texttt{get\_data\_mover()}:
    set/return the workflow default data mover.
  \item \texttt{set\_stager\_mover(stager\_mover)}: set the workflow default
    staging mode.
  \item \texttt{get\_scratch\_dir\_base()}: compute or return the workflow
    scratch base.
  \item \texttt{find\_task\_by\_name(workflow\_name, task\_name)}: locate a
    task.
  \item \texttt{as\_json()}: serialize workflow metadata.
  \item \texttt{load\_json(Json\_data)}: load tasks from JSON-like data.
  \item \texttt{Validate\_WF()}: reject cycles.
\end{itemize}

Two of these methods deserve a cross-reference a pure API listing would
otherwise omit. \texttt{as\_json()}/\texttt{load\_json()} are the mechanism
underlying the cross-process workflow transport implied by
\Cref{sec:service-architecture}'s optional DAGon service, and are also what
makes a serialized workflow --- including its \texttt{LLM} task provider
references, per \Cref{ch:llm-tasks} --- portable across machines, provided
the receiving machine has the same named provider configured locally.
\texttt{get\_scratch\_dir\_base()} is the method whose return value
ultimately determines where every artifact discussed in
\Cref{ch:lessons-03-04}'s scratch-directory anatomy (\Cref{fig:task-dir})
physically lives.

\section{\texttt{dagon.Status}}

\begin{center}
\begin{tabular}{@{}l@{}}
\toprule
\textbf{Value} \\
\midrule
\texttt{Status.READY} \\
\texttt{Status.WAITING} \\
\texttt{Status.RUNNING} \\
\texttt{Status.FINISHED} \\
\texttt{Status.FAILED} \\
\bottomrule
\end{tabular}
\end{center}

This five-state model is the concrete realization of the abstract,
topological-order-driven scheduling argument developed in
\Cref{sec:topological-ordering}: \texttt{READY} corresponds exactly to a
node with no remaining incoming edges in the source-removal construction
given there, and the transition from \texttt{WAITING} to \texttt{READY},
triggered when a task's last predecessor reaches \texttt{FINISHED}, is the
runtime enactment of that construction's repeated source-removal step.

\section{\texttt{dagon.task.TaskType}}

\begin{center}
\begin{tabular}{@{}l@{}}
\toprule
\textbf{Value} \\
\midrule
\texttt{TaskType.CHECKPOINT} \\
\texttt{TaskType.BATCH} \\
\texttt{TaskType.SLURM} \\
\texttt{TaskType.CLOUD} \\
\texttt{TaskType.DOCKER} \\
\texttt{TaskType.LLM} \\
\texttt{TaskType.NATIVE} \\
\texttt{TaskType.WEB} \\
\bottomrule
\end{tabular}
\end{center}

\section{\texttt{dagon.task.DagonTask}}

\begin{lstlisting}
task = DagonTask(TaskType.BATCH, "A", "echo A")
\end{lstlisting}

The concrete class is selected from the \texttt{tasks\_types} mapping:

\begin{center}
\begin{tabular}{@{}lll@{}}
\toprule
\textbf{Task type} & \textbf{Module} & \textbf{Class} \\
\midrule
\texttt{CHECKPOINT} & \texttt{dagon.checkpoint} & \texttt{Checkpoint} \\
\texttt{BATCH}       & \texttt{dagon.batch} & \texttt{Batch} \\
\texttt{CLOUD}       & \texttt{dagon.remote} & \texttt{CloudTask} \\
\texttt{DOCKER}      & \texttt{dagon.docker\_task} & \texttt{DockerTask} \\
\texttt{LLM}         & \texttt{dagon.llm} & \texttt{LLMTask} \\
\texttt{NATIVE}      & \texttt{dagon.native} & \texttt{NativeTask} \\
\texttt{WEB}         & \texttt{dagon.web} & \texttt{WebTask} \\
\texttt{SLURM}       & \texttt{dagon.batch} & \texttt{Slurm} \\
\bottomrule
\end{tabular}
\end{center}

This table is, in effect, the implementation-level evidence for the factory-
pattern discussion of \Cref{sec:factory-pattern-depth}: every row maps an
identical calling convention --- \texttt{DagonTask(TaskType.X, name,
payload, ...)} --- onto a distinct concrete class living in its own module,
several of which (\texttt{dagon.docker\_task}, \texttt{dagon.llm},
\texttt{dagon.native}, \texttt{dagon.web}) are, per
\Cref{sec:known-constraints}, lazily imported, so that constructing a
\texttt{BATCH} task never requires the Docker SDK or any web-request library
to even be installed.

\section{\texttt{dagon.task.Task}}

\begin{lstlisting}[caption={Base constructor.}]
Task(name, command, working_dir=None, transversal_workflow=None, globusendpoint=None)
\end{lstlisting}

\paragraph{Important attributes.} \texttt{name}, \texttt{command},
\texttt{working\_dir}, \texttt{nexts}, \texttt{prevs},
\texttt{reference\_count}, \texttt{workflow}, \texttt{status},
\texttt{data\_mover}, \texttt{stager\_mover}, \texttt{globusendpoint}.

\paragraph{Important methods.} \texttt{add\_dependency\_to(task)},
\texttt{add\_transversal\_point(task)}, \texttt{pre\_run()},
\texttt{pre\_process\_command(command)}, \texttt{post\_process\_command(command)},
\texttt{include\_command(body)}, \texttt{create\_working\_dir()},
\texttt{execute()}, \texttt{run()}, \texttt{as\_json()},
\texttt{get\_scratch\_dir()}, \texttt{get\_endpoint()},
\texttt{set\_endpoint(globusendpoint)}, \texttt{set\_data\_mover(data\_mover)},
\texttt{set\_stager\_mover(stager\_mover)}.

Subclasses are expected to implement or specialize command execution
(\Cref{ch:architecture}). Of the attributes listed, \texttt{prevs} and
\texttt{nexts} are the two a workflow author will inspect most often during
the local-structural-verification pattern recommended throughout
\Cref{ch:backends,ch:lessons-01-02}; \texttt{pre\_run()} and
\texttt{pre\_process\_command()} are the two methods that jointly implement
the dependency-inference algorithm formalized in \Cref{ch:schema}.

\section{Batch and Slurm Tasks}

\begin{lstlisting}
DagonTask(TaskType.BATCH, "local", "hostname")
\end{lstlisting}

Remote batch is selected by passing SSH parameters:

\begin{lstlisting}
DagonTask(
    TaskType.BATCH, "remote", "hostname",
    ip="example.org", ssh_username="user", keypath="/path/to/key",
)
\end{lstlisting}

\begin{lstlisting}
DagonTask(
    TaskType.SLURM, "simulation", "./run.sh",
    partition="compute", ntasks=4, memory=4096, time="01:00:00",
)
\end{lstlisting}

Supported Slurm constructor options include: \texttt{comment},
\texttt{partition}, \texttt{ntasks}, \texttt{memory}, \texttt{time},
\texttt{nodes}, \texttt{ntasks\_per\_node}, \texttt{working\_dir},
\texttt{globusendpoint}. Remote Slurm is selected by passing \texttt{ip}.
See \Cref{ch:backends} for the queueing-cost and structural-verification
discussion that should accompany any first use of these parameters against
a real cluster.

\section{Docker Tasks}

\begin{lstlisting}
DagonTask(
    TaskType.DOCKER, "container-task", "python --version",
    image="python:3.12-slim",
)
\end{lstlisting}

Options include: \texttt{image}, \texttt{container\_id},
\texttt{working\_dir}, \texttt{globusendpoint}, \texttt{remove},
\texttt{volume}. Remote Docker is selected by passing SSH information
including \texttt{ip}. The \texttt{image} value, as discussed in
\Cref{ch:backends}'s tag-versus-digest reproducibility comparison, should be
considered for digest-pinning whenever the workflow's results are intended
for archival or publication.

\section{Cloud Tasks}

Cloud tasks are backed by \texttt{dagon.remote.CloudTask} and Apache
Libcloud. They require provider information, an SSH user, key options, and
either an existing instance or instance flavor data. Cloud behavior is
provider-specific and should be tested in the target environment. As noted
in \Cref{ch:configuration}'s discussion of configuration sections not
exercised in this book's worked examples, no runnable cloud example is
given here; the provider-specific authentication and provisioning model
must be verified against the target provider's own current documentation.

\section{LLM Tasks}

\begin{lstlisting}[caption={\texttt{LLMTask} constructor.}]
DagonTask(TaskType.LLM, name, prompt, provider, params=None, input_files=None,
          working_dir=None, output_file="response.json", timeout=120)
\end{lstlisting}

\texttt{prompt} is a JSON object (or a JSON string) containing
\texttt{messages}; string values may use \texttt{\{parameter\}} placeholders
from \texttt{params} or \texttt{input\_files}. Provider configuration is read
from \texttt{[llm.<provider>]}. \texttt{input\_files} maps parameter names to
\texttt{workflow://} references, which infer dependencies and stage local
UTF-8 text before the request. The default \texttt{timeout=120} should be
reconsidered, per \Cref{ch:llm-tasks}'s discussion of operational
constraints, against the documented typical latency of whichever real
provider a workflow ultimately targets. See \Cref{ch:llm-tasks}.

\section{Native Tasks}

\begin{lstlisting}[caption={\texttt{NativeTask} constructor.}]
DagonTask(TaskType.NATIVE, name, function, inputs=None, outputs=None,
          executor="local", resources=None, python=sys.executable,
          working_dir=None, environment=None)
\end{lstlisting}

Inputs may be JSON scalars, existing local files, or \texttt{workflow://}
file references. File values are staged below \texttt{inputs/};
\texttt{outputs} maps parameter names to safe relative paths below
\texttt{outputs/}. JSON return metadata is written to
\texttt{.dagon/native\_result.json}. \texttt{executor} accepts
\texttt{local} and \texttt{slurm}; functions must be importable module-level
callables, for the serialization-related reasons developed in full in
\Cref{ch:native-tasks}. See \Cref{ch:native-tasks}.

\section{Web Tasks}

\begin{lstlisting}[caption={\texttt{WebTask} specification.}]
DagonTask(TaskType.WEB, name,
    {"method": "GET", "url": "https://example.org",
     "outputs": {"body": "response.json"}},
    executor="local")
\end{lstlisting}

Specifications support \texttt{query}, \texttt{headers}, \texttt{json},
\texttt{data}, \texttt{multipart}, \texttt{auth}, \texttt{timeout}, retry
settings, expected status codes, and \texttt{outputs}. Recursive
\texttt{workflow://} values infer dependencies; referenced files are copied
below \texttt{inputs/}. Body, headers, and metadata outputs are safe
relative paths below \texttt{outputs/}; \texttt{.dagon/web\_result.json}
stores non-secret status metadata. Supported authentication is bearer,
basic, API-key header, or API-key query, always through environment-variable
\emph{names} rather than literal secrets. \texttt{executor="slurm"} accepts
existing scheduler options in \texttt{resources}. The retry settings
mentioned here interact directly with the idempotency considerations
developed in \Cref{ch:web-tasks}; review that chapter's discussion before
enabling automatic retries against a non-idempotent endpoint. See
\Cref{ch:web-tasks}.

\section{Data Movement Enumerations}

\paragraph{\texttt{DataMover}.} \texttt{DONTMOVE}, \texttt{LINK},
\texttt{COPY}, \texttt{SCP}, \texttt{HTTP}, \texttt{HTTPS}, \texttt{FTP},
\texttt{SFTP}, \texttt{GRIDFTP}, \texttt{SKYCDS}. The cost/self-containment
trade-off between \texttt{COPY} and \texttt{LINK} specifically is examined
quantitatively in \Cref{ch:lessons-05-07}.

\paragraph{\texttt{StagerMover}.} \texttt{NORMAL}, \texttt{PARALLEL},
\texttt{SLURM}.

\section{Configuration API}

\begin{lstlisting}
from dagon.config import read_config

all_sections = read_config("dagon.ini")
batch = read_config("dagon.ini", "batch")
\end{lstlisting}

If a requested section does not exist, \texttt{read\_config(..., section)}
returns \texttt{None}. This function underlies the file-based configuration
path discussed in \Cref{ch:configuration}; note that it is bypassed entirely
when a \texttt{Workflow} is constructed with an in-memory \texttt{config}
dictionary instead.

\section{Workflow Schema}

The \texttt{workflow://} schema is documented in full in \Cref{ch:schema}.
In brief:

\begin{lstlisting}
workflow://<workflow-name>/<task-name>/<path-inside-task-working-directory>
workflow:///<task-name>/<path-inside-task-working-directory>
\end{lstlisting}

\section{Meta-Workflows}

\begin{lstlisting}
from dagon.dag_tps import DAG_TPS

meta = DAG_TPS("ExperimentSet")
meta.add_workflow(workflow_a)
meta.add_workflow(workflow_b)
meta.make_dependencies()
meta.run()
\end{lstlisting}

Useful for transversal workflows that reference data across workflow
boundaries (\Cref{ch:meta-async}). The meta-workflow code is
legacy-sensitive and less covered by automated tests than local
\texttt{Workflow} behavior; production use should be verified with
representative workflows, preserving the verification command with the
experiment record. See also this appendix's framing of this caveat within
\Cref{ch:synthesis}'s closing, honesty-focused revisiting of the book's
central claims.

\section{Service API}

\texttt{dagon.api.API} communicates with a DAGon service endpoint. It
supports: \texttt{create\_workflow(workflow)},
\texttt{add\_task(workflow\_id, task)},
\texttt{update\_task\_status(workflow\_id, task, status)},
\texttt{get\_task(workflow\_id, task)},
\texttt{get\_workflow\_by\_name(workflow\_name)},
\texttt{update\_task(workflow\_id, task, attribute, value)},
\texttt{add\_dependency(workflow\_id, task, dependency)}.

The service is optional and disabled in the default sample configuration
(\Cref{sec:service-architecture}). A reader who has worked through this
book's local-first examples exclusively, as the examples-catalog discussion
of \Cref{ch:examples-catalog} recommends, will not have exercised this API
directly; it becomes relevant only when a deployment genuinely requires the
multi-process or multi-host coordination discussed in
\Cref{ch:architecture}'s comparison of layered-in-process and microservice
architectural alternatives.
